\documentclass[12pt,letterpaper, onecolumn]{exam}
\usepackage{bm}
\usepackage{booktabs}
\usepackage{amsmath,amsfonts}
\usepackage{amssymb,amsthm}
\usepackage[lmargin=71pt, tmargin=1.2in]{geometry}  %For centering solution box
\geometry{top=25.4mm,bottom=25.4mm,left=20mm,right=20mm,headheight=2.17cm,headsep=4mm,footskip=12mm}
% \lhead{Leaft Header\\}
% \rhead{Right Header\\}
% \chead{\hline} % Un-comment to draw line below header
\thispagestyle{empty}   %For removing header/footer from page 1

\begin{document}

\begingroup  
    \centering
    \LARGE Machine Learning\\
    \LARGE Homework 4\\[0.5em]
    \large Dec. 2025\\[0.5em]
    \large Fill your Name\par
    \large Fill your Student ID\par
\endgroup
\rule{\textwidth}{0.4pt}
\pointsdroppedatright   %Self-explanatory
\printanswers
\renewcommand{\solutiontitle}{\noindent\textbf{Ans:}\enspace}   %Replace "Ans:" with starting keyword in solution box

\begin{questions} 
    \question Here is a sample table. Please use the data in the table~\ref{tab:decision_tree} to solve the decision tree-related problems.
    \begin{table}[!h]
        \centering
        \begin{tabular}{c|ccccc}
        \toprule[2pt]
            \textbf{Day} & \textbf{Outlook} & \textbf{Temperature} & \textbf{Humidity} & \textbf{wind} & \textbf{Play Tennis}\\
            \midrule
            D1& Sunny & Hot & High & Weak & No \\
            D2& Sunny & Hot & High & Strong & No \\
            D3& Overcast & Hot & High & Weak & Yes \\
            D4& Rain & Mild & High & Weak & Yes \\
            D5& Rain & Cool & Normal & Weak & Yes \\
            D6& Rain & Cool & Normal & Strong & No \\
            D7& Overcast & Cool & Normal & Strong & Yes \\
            D8& Sunny & Mild & High & Weak & No \\
            D9& Sunny & Cool & Normal & Weak & Yes \\
            D10& Rain & Mild & Normal & Weak & Yes \\
            D11& Sunny & Mild & Normal & Strong & Yes \\
            D12& Overcast & Mild & High & Strong & Yes \\
            D13& Overcast & Hot & Normal & Weak & Yes \\
            D14& Rain & Hot & High & Strong & No \\
            
        \bottomrule[2pt]
        \end{tabular}
        \caption{Instances for building a decision tree.}
        \label{tab:decision_tree}
    \end{table}
    
    \begin{parts}
        \part Compute the Overall Entropy.
        \part Compute the information gain and select the root node of the decision tree.
    \end{parts}
    
    \begin{solution}
        \begin{parts}
            \part 
            
            \part
            
        \end{parts}
    \end{solution}

    \newpage
    \question Here is a sample table. Please use the data in the table~\ref{tab:naive_bayes} to solve the naive Bayes-related problems.
        \begin{table}[!h]
            \centering
            
            \begin{tabular}{c|ccccc}
            \toprule[2pt]
                \textbf{ID} & \textbf{age} & \textbf{income} & \textbf{student} & \textbf{credit rating} & \textbf{buy computer}\\
                \midrule
                1 & \textless30 & high & no & fair & no\\
                2 & \textless30 & high & no & excellent & no\\
                3 & 30-40 & high & no & fair & yes\\
                4 & \textgreater40 & medium & no & fair & yes\\
                5 & \textgreater40 & low & yes & fair & yes\\
                6 & \textgreater40 & low & yes & excellent & no\\
                7 & 30-40 & low & yes & excellent & yes\\
                8 & \textless30 & medium & no & fair & no\\
                9 & \textless30 & low & yes & fair & yes\\
                10 & \textgreater40 & medium & yes & fair & yes\\
                11 & \textless30 & medium & yes & excellent & yes\\
                12 & 30-40 & medium & no & excellent & yes\\
                13 & 30-40 & high & yes & fair & yes\\
                14 & \textgreater40 & medium & no & excellent & no\\
            \bottomrule[2pt]
            \end{tabular}
            \caption{Instances for naive Bayes}
            \label{tab:naive_bayes}
        \end{table}

    \begin{parts}
        \part Computer the prior probability.
        \part Computer the conditional probability of "age".
        \part Given the new samples: x = (age$<$30, income=medium, student=yes, credit rating=fair). Based on Naive Bayesian calculations, it determines whether to purchase a computer.
    \end{parts}
    
    
    \begin{solution}
        \begin{parts}
            \part 

            \part 

            \part
           
        \end{parts}
    \end{solution}

    \newpage

    \question Perform clustering analysis on the following two-dimensional data using the k-medoids algorithm. Assuming k=2, initialize the cluster centers using data points G and H. Calculate the similarity between any two points using Manhattan distance.

    \begin{table}[!h]
        \centering
        \begin{tabular}{c|cccccccc}
            \toprule[2pt]
             ID       & A & B & C & D & E & F & G & H \\
             \midrule
             feature 1& 1.2 & 2.1 & 1.9 & 2.3 & 3.1 & 4.2 & 1.2 & 3.3 \\
             feature 2& 1.1 & 2.4 & 3.1 & 4.2 & 4 & 4 & 1 & 2.2 \\
             \bottomrule[2pt]
        \end{tabular}
        \caption{Instances for the k-medoids}
        \label{tab:cluster}
    \end{table}
    
    \begin{parts}
        \part  Provide the clustering results and key computational steps.
    \end{parts}


    \begin{solution}
        \begin{parts}
            \part 
            
        \end{parts} 
    \end{solution}

    
\end{questions}
\end{document}


